\documentclass[12pt,a4paper]{article}

% -------------------------------------------------
% Encoding and language
% -------------------------------------------------
\usepackage[utf8]{inputenc}
\usepackage[T1]{fontenc}
\usepackage[spanish,es-nodecimaldot]{babel}

% -------------------------------------------------
% Page layout
% -------------------------------------------------
\usepackage[
a4paper,
top=2.2cm,
bottom=2.2cm,
left=2.5cm,
right=2.5cm
]{geometry}

% -------------------------------------------------
% Typography
% -------------------------------------------------
\usepackage{lmodern}
\usepackage{microtype}

% -------------------------------------------------
% Mathematics
% -------------------------------------------------
\usepackage{amsmath}
\usepackage{amssymb}

% -------------------------------------------------
% Tables and lists
% -------------------------------------------------
\usepackage{array}
\usepackage{longtable}
\usepackage{booktabs}
\usepackage{enumitem}

% -------------------------------------------------
% Colors
% -------------------------------------------------
\usepackage{xcolor}

\definecolor{primary}{RGB}{31,78,121}
\definecolor{secondary}{RGB}{68,68,68}
\definecolor{lightgray}{RGB}{245,245,245}

% -------------------------------------------------
% Headers and footers
% -------------------------------------------------
\usepackage{fancyhdr}

\pagestyle{fancy}
\fancyhf{}

\fancyhead[L]{\small Planeación Periódica de la Enseñanza}
\fancyhead[R]{\small PGF-02-R39}

\fancyfoot[C]{\thepage}

\renewcommand{\headrulewidth}{0.4pt}
\renewcommand{\footrulewidth}{0pt}

% -------------------------------------------------
% Section formatting
% -------------------------------------------------
\usepackage{titlesec}

\titleformat{\section}
{\Large\bfseries\color{primary}}
{\thesection.}
{0.5em}
{}

\titleformat{\subsection}
{\large\bfseries\color{primary}}
{}
{0pt}
{}

\titleformat{\subsubsection}
{\normalsize\bfseries\color{secondary}}
{}
{0pt}
{}

% -------------------------------------------------
% Links
% -------------------------------------------------
\usepackage[
colorlinks=true,
linkcolor=primary,
urlcolor=primary,
citecolor=primary
]{hyperref}

% -------------------------------------------------
% Paragraph formatting
% -------------------------------------------------
\setlength{\parindent}{0pt}
\setlength{\parskip}{0.6em}

% -------------------------------------------------
% List formatting
% -------------------------------------------------
\setlist[itemize]{
	leftmargin=1.5cm,
	itemsep=0.2em,
	topsep=0.2em
}

\setlist[description]{
	leftmargin=1.5cm,
	labelindent=0cm,
	style=nextline,
	itemsep=0.5em
}

% -------------------------------------------------
% Custom commands
% -------------------------------------------------

\newcommand{\criterio}[2]{
	\item[\textbf{#1:}] #2
}

\newcommand{\placeholder}[1]{
	\textit{\textcolor{secondary}{[#1]}}
}

% -------------------------------------------------
% Document
% -------------------------------------------------

\begin{document}
	
	% =================================================
	% TITLE
	% =================================================
	
	\begin{center}
		
		{\Large\bfseries\color{primary}
			PLANEACIÓN PERIÓDICA DE LA ENSEÑANZA}
		
		\vspace{0.4cm}
		
		{\normalsize\bfseries PGF-02-R39}
		
		\vspace{0.3cm}
		
		{\normalsize
			\textbf{Fecha:} 12-08-2026 a 05-11-2026}
		
	\end{center}
	
	\vspace{0.8cm}
	
	% =================================================
	% 1. IDENTIFICACIÓN
	% =================================================
	
	\section{Identificación}
	
	\textbf{Nombre del ámbito / habilidad:}
	
	Seminarios -- ``Emprendimiento e innovación''
	
	\textbf{Pregunta de investigación del seminario:}
	
	¿Cómo puede la inteligencia artificial ayudarnos a diseñar y construir soluciones tecnológicas para problemas de nuestro entorno?
	
	\textbf{Nombre del profesor:}
	
	Nicolás López
	
	\textbf{Grado:}
	
	10°
	
	\textbf{Periodo:}
	
	37
	
	% =================================================
	% 2. OBJETIVO DE GRADO
	% =================================================
	
	\section{Objetivo de grado}
	
	\placeholder{Información no especificada en el documento proporcionado.}
	
	% =================================================
	% 3. OBJETIVO DE APRENDIZAJE
	% =================================================
	
	\section{Objetivo de aprendizaje}
	
	\subsection*{Objetivo de aprendizaje}
	
	Identificar necesidades, problemas u oportunidades de innovación en el entorno cercano mediante observación contextual, empatía y definición de retos.
	
	\subsection*{Criterios de evaluación}
	
	\begin{description}
		
		\criterio{C1}{
			Reconoce necesidades, incomodidades u oportunidades de mejora
			en la comunidad escolar o cercana.
		}
		
		\criterio{C2}{
			Observa comportamientos, espacios, usuarios o situaciones para
			comprender mejor una necesidad.
		}
		
		\criterio{C3}{
			Escucha a posibles usuarios mediante conversaciones, preguntas
			orientadas o ejercicios de empatía.
		}
		
		\criterio{C4}{
			Diferencia una necesidad real de una idea de solución anticipada.
		}
		
		\criterio{C5}{
			Organiza información del contexto en mapas, perfiles de usuario,
			listados de hallazgos o diagramas simples.
		}
		
		\criterio{C6}{
			Prioriza una oportunidad de innovación considerando relevancia,
			posibilidad de acción y beneficio esperado.
		}
		
		\criterio{C7}{
			Define un reto de innovación claro, concreto y comprensible
			para el grupo.
		}
		
		\criterio{C8}{
			Relaciona el reto con saberes de diseño, comunicación, economía,
			tecnología, ciudadanía u otra disciplina pertinente.
		}
		
		\criterio{C9}{
			Participa en dinámicas de ideación inicial con apertura a
			propuestas diversas.
		}
		
		\criterio{C10}{
			Valora la importancia de comprender al usuario antes de
			proponer una solución.
		}
		
	\end{description}
	
	% =================================================
	% 4. ESCENARIO Y NARRATIVA
	% =================================================
	
	\section{Escenario y narrativa}
	
	\textbf{Escenario:}
	
	Es el contexto o marco general en el que se ubican las experiencias
	de aprendizaje propuestas para el periodo. Incluye el lugar, tiempo,
	circunstancias, problemáticas y demás elementos relevantes.
	
	\textbf{Narrativa:}
	
	Descripción o construcción dialógica realizada por el equipo de grado
	para alinear las acciones pedagógicas.
	
	\placeholder{Información específica del escenario y la narrativa no especificada
		en el documento proporcionado.}
	
	% =================================================
	% 5. REFERENTES CONCEPTUALES
	% =================================================
	
	\section{Referentes conceptuales}
	
	\textbf{Innovación:}
	
	Según la Organización para la Cooperación y el Desarrollo Económicos
	(OCDE) y Eurostat (2018), la innovación es el proceso mediante el cual
	se crean o mejoran productos, servicios, procesos o formas de
	organización para responder de manera más eficiente, creativa y
	sostenible a una necesidad o problema, generando valor para las
	personas y la sociedad.
	
	La innovación generalmente:
	
	\begin{itemize}
		
		\item Responde a una necesidad o problema concreto.
		
		\item Aporta una mejora significativa frente a soluciones existentes.
		
		\item Genera valor para las personas o la sociedad.
		
		\item Puede ser tecnológica, social, educativa, ambiental u
		organizacional.
		
		\item Requiere creatividad, conocimiento y capacidad para
		implementar ideas.
		
	\end{itemize}
	
	% =================================================
	% 6. PRE-LECCIÓN Y CONTEXTO
	% =================================================
	
	\section{Pre-lección y contexto}
	
	Se realizará un encuentro con toda la promoción en Sede Social para
	explicar a los estudiantes la proyección y el alcance del nuevo
	formato de seminarios.
	
	Durante este encuentro:
	
	\begin{itemize}
		
		\item Se presentará el equipo docente a cargo.
		
		\item Se proyectarán las preguntas de interés de los diferentes
		seminarios.
		
		\item Se invitará a los estudiantes a ingresar a sus correos en
		la tarde, al llegar a casa, para realizar la inscripción.
		
		\item Se explicarán los aspectos logísticos relacionados con el
		desarrollo de los seminarios.
		
	\end{itemize}
	
	Posteriormente, se contará con un invitado especial de la Universidad
	CESA, quien explicará a los estudiantes qué significa emprender, cómo
	hacerlo y cuál es su importancia para la vida adulta.
	
	Se promoverá la participación de los estudiantes mediante preguntas,
	opiniones y reflexiones relacionadas con el emprendimiento y la
	innovación.
	
	% =================================================
	% 7. EXPERIENCIA
	% =================================================
	
	\section{Experiencia}
	
	Las experiencias corresponden a las acciones pedagógicas que permiten
	desarrollar el objetivo de aprendizaje, atendiendo a los criterios de
	evaluación. Se deben incluir los enlaces a los recursos virtuales o
	las rutas utilizadas para compartir las experiencias (OPEN, Teams,
	etc.).
	
	% -------------------------------------------------
	% Acción pedagógica 1
	% -------------------------------------------------
	
	\subsection{Acción pedagógica N.\textsuperscript{o} 1}
	
	\textbf{Criterios de evaluación:} C1 a C4
	
	\subsubsection*{Laboratorio de Innovación:
		Identificando oportunidades en mi comunidad}
	
	Organizados en equipos de trabajo, los estudiantes realizarán un
	proceso de observación e investigación sobre una problemática que
	afecte al colegio.
	
	Algunos posibles temas son:
	
	\begin{itemize}
		
		\item Manejo de residuos.
		
		\item Movilidad.
		
		\item Convivencia.
		
		\item Salud mental.
		
		\item Uso responsable de la tecnología.
		
		\item Emprendimiento local.
		
		\item Seguridad.
		
		\item Acceso a espacios recreativos.
		
		\item Otras problemáticas identificadas por los estudiantes.
		
	\end{itemize}
	
	Para ello, aplicarán diferentes estrategias de recolección de
	información, como:
	
	\begin{itemize}
		
		\item Encuestas.
		
		\item Entrevistas.
		
		\item Observaciones.
		
		\item Búsqueda y análisis de información.
		
	\end{itemize}
	
	Posteriormente, analizarán las causas, consecuencias e impacto del
	problema y presentarán un diagnóstico sustentado con evidencias.
	
	Para esta acción pedagógica, cada equipo contará con cuatro clases.
	El maestro realizará un seguimiento continuo de los avances y
	procesos desarrollados por los estudiantes.
	
	% -------------------------------------------------
	% Acción pedagógica 2
	% -------------------------------------------------
	
	\subsection{Acción pedagógica N.\textsuperscript{o} 2}
	
	\textbf{Criterios de evaluación:} C5 a C10
	
	A partir del diagnóstico realizado previamente, los equipos
	desarrollarán una propuesta de solución innovadora que sea viable,
	sostenible y genere un impacto positivo en la comunidad.
	
	La propuesta podrá materializarse como:
	
	\begin{itemize}
		
		\item Producto.
		
		\item Servicio.
		
		\item Campaña.
		
		\item Aplicación.
		
		\item Estrategia social.
		
		\item Iniciativa ambiental.
		
		\item Proyecto comunitario.
		
	\end{itemize}
	
	Durante el proceso, los estudiantes deberán:
	
	\begin{itemize}
		
		\item Justificar su idea.
		
		\item Identificar los recursos necesarios.
		
		\item Identificar los posibles beneficiarios.
		
		\item Definir cómo evaluarán el impacto de su solución.
		
	\end{itemize}
	
	Cada área del conocimiento podrá aportar elementos propios de su
	disciplina para fortalecer las propuestas, tales como:
	
	\begin{itemize}
		
		\item Diseño del modelo de negocio.
		
		\item Elaboración de presupuestos.
		
		\item Análisis científico.
		
		\item Comunicación efectiva.
		
		\item Identidad visual.
		
		\item Prototipado.
		
		\item Sostenibilidad.
		
		\item Habilidades ciudadanas.
		
		\item Presentación y comunicación del proyecto.
		
	\end{itemize}
	
	% =================================================
	% 7.1 PAUSA
	% =================================================
	
	\subsection{7.1 Pausa: ¿Qué voy logrando?}
	
	\textbf{Descripción:}
	
	Descripción en torno a ``¿Qué voy logrando?''
	
	\vspace{0.3cm}
	
	\textbf{Modo de proceder:}
	
	\placeholder{Por definir.}
	
	\textbf{Modo de ser:}
	
	\placeholder{Por definir.}
	
	% =================================================
	% 7.2 PAUSA
	% =================================================
	
	\subsection{7.2 Pausa: ¿Qué logré?}
	
	\textbf{Descripción:}
	
	Descripción en torno a ``¿Qué logré?''
	
	\vspace{0.3cm}
	
	\textbf{Modo de proceder:}
	
	\placeholder{Por definir.}
	
	\textbf{Modo de ser:}
	
	\placeholder{Por definir.}
	
	% =================================================
	% REFERENCES
	% =================================================
	
	\section*{Referencias}
	
	\placeholder{Por completar.}
	
\end{document}